\input{../../../templates/es/preambulo.tex}

% Los listados se toman de src/ con \lstinputlisting; la diapositiva es el archivo.
\lstset{
    basicstyle=\ttfamily\fontsize{8pt}{9.6pt}\selectfont,
    breaklines=true,
    breakatwhitespace=true,
    columns=fullflexible,
    keepspaces=true,
    xleftmargin=0.3em,
    framesep=2pt,
    aboveskip=0pt,
    belowskip=0pt
}

\newcommand{\marcoresultado}[3]{%
    \begin{frame}{#1}
        \begin{center}
            \includegraphics[height=0.62\textheight,width=\textwidth,keepaspectratio]{#2}
        \end{center}
        \vspace{0.25em}
        {\small #3\par}
    \end{frame}%
}

\newcommand{\marcoafirmacion}[2]{%
    \begin{frame}{#1}
        {\small #2\par}
    \end{frame}%
}

\date{}

\begin{document}

\tituloleccion{01}{Tu primer algoritmo genético}

\section{Esta lección}

\begin{frame}{Esta lección}
El encuadre mostró cómo se lee una lección y cómo se corre el código. Hoy se construye el primer algoritmo, un operador a la vez.

\vspace{0.6em}
\begin{itemize}
    \item Una población, no un punto; la selección descarta, la cruza recombina, la mutación inventa.
    \item El ciclo de esos tres, unas cuantas generaciones.
    \item El mismo código, otra semilla, una colina local: un AG no es una garantía.
\end{itemize}
\end{frame}

% ================= primer_ejemplo_01_el_paisaje =================
\begin{frame}{Modelo matemático}
\[
\max_{x\in[-10,10]} f(x),\qquad f(x)=\sin x-0.2|x|
\]
Un individuo guarda un candidato \(x\) y su aptitud \(f(x)\). Una población es
una muestra finita, así que su campeón es el \emph{mejor observado}. La
cuadrícula densa es una referencia numérica, no una prueba del maximizador.
\end{frame}

\section{El problema}

\begin{frame}{El problema}
Antes de escribir cualquier algoritmo, veamos qué estamos buscando.
\end{frame}

\marcoresultado{El problema}{../figuras/primer_ejemplo_01_el_paisaje.png}{%
    Hay cuatro cimas; subir por la pendiente se detiene en aquella donde empezó}

% ================= primer_ejemplo_02_poblacion_aleatoria =================
\section{Una población}

\begin{frame}{Una población}
Un algoritmo genético no sigue la pista de una sola solución candidata, sigue
muchas. Aquí simplemente esparcimos diez al azar y vemos dónde caen.
\end{frame}

\begin{frame}[fragile]{(1a) Individuo: estado y aptitud}
\lstinputlisting[basicstyle=\ttfamily\fontsize{7pt}{8.5pt}\selectfont,firstline=46,lastline=62]{../src/primer_ejemplo_02_poblacion_aleatoria.py}
\end{frame}

\begin{frame}[fragile]{(1b) Individuo: acceso y representación}
\lstinputlisting[basicstyle=\ttfamily\fontsize{7pt}{8.5pt}\selectfont,firstline=64,lastline=71]{../src/primer_ejemplo_02_poblacion_aleatoria.py}
\end{frame}

\begin{frame}[fragile]{(2) crear\_aleatorio()}
\lstinputlisting[basicstyle=\ttfamily\fontsize{7pt}{8.5pt}\selectfont,firstline=74,lastline=88]{../src/primer_ejemplo_02_poblacion_aleatoria.py}
\end{frame}

\begin{frame}[fragile]{(3) la población}
\lstinputlisting[basicstyle=\ttfamily\fontsize{7pt}{8.5pt}\selectfont,firstline=91,lastline=114]{../src/primer_ejemplo_02_poblacion_aleatoria.py}
\end{frame}

\marcoresultado{Una población}{../figuras/primer_ejemplo_02_poblacion_aleatoria.png}{%
    Diez suposiciones a ciegas no caen cerca del óptimo}

% ================= primer_ejemplo_03_seleccion =================
\section{Selección}

\begin{frame}{Selección}
La población del paso 2 es aleatoria y en su mayoría mala. La selección es la
primera fuerza de la evolución: decide quién llega a reproducirse. Aquí no se
crea nada nuevo: solo elegimos, con repetición, entre lo que ya existe.
\end{frame}

\begin{frame}[fragile]{(1) seleccion\_torneo()}
\lstinputlisting[basicstyle=\ttfamily\fontsize{6.5pt}{7.8pt}\selectfont,firstline=79,lastline=103]{../src/primer_ejemplo_03_seleccion.py}
\end{frame}

\marcoresultado{Selección}{../figuras/primer_ejemplo_03_seleccion.png}{%
    No aparece ningún valor nuevo de \texttt{x} — la selección solo copia; 4 de 10 individuos se extinguen}

% ================= primer_ejemplo_04_cruza =================
\section{Cruza}

\begin{frame}{Cruza}
La selección por sí sola solo puede copiar. La cruza es el primer operador que
crea valores que antes no estaban en la población, combinando dos progenitores.
\end{frame}

\begin{frame}[fragile]{(1) acotar()}
\lstinputlisting[basicstyle=\ttfamily\fontsize{7pt}{8.5pt}\selectfont,firstline=46,lastline=65]{../src/primer_ejemplo_04_cruza.py}
\end{frame}

\begin{frame}[fragile]{(2) cruza\_mezcla()}
\lstinputlisting[basicstyle=\ttfamily\fontsize{6pt}{7.2pt}\selectfont,firstline=88,lastline=118]{../src/primer_ejemplo_04_cruza.py}
\end{frame}

\begin{frame}[fragile]{(3) cruza()}
\lstinputlisting[basicstyle=\ttfamily\fontsize{7pt}{8.5pt}\selectfont,firstline=121,lastline=137]{../src/primer_ejemplo_04_cruza.py}
\end{frame}

\marcoresultado{Cruza}{../figuras/primer_ejemplo_04_cruza.png}{%
    14 de 20 hijos caen fuera del intervalo de los progenitores — eso es lo que compra \texttt{alfa > 0}}

% ================= primer_ejemplo_05_mutacion =================
\section{Mutación}

\begin{frame}{Mutación}
Cuando los genes de los progenitores quedan muy agrupados en una colina, esta
cruza tiene poco alcance. La mutación agrega un desplazamiento independiente y
todavía puede proponer valores lejanos.
\end{frame}

\begin{frame}[fragile]{(1) mutacion\_gaussiana()}
\lstinputlisting[basicstyle=\ttfamily\fontsize{6.5pt}{7.8pt}\selectfont,firstline=82,lastline=103]{../src/primer_ejemplo_05_mutacion.py}
\end{frame}

\begin{frame}[fragile]{(2) mutar()}
\lstinputlisting[basicstyle=\ttfamily\fontsize{7pt}{8.5pt}\selectfont,firstline=105,lastline=119]{../src/primer_ejemplo_05_mutacion.py}
\end{frame}

\marcoresultado{Mutación}{../figuras/primer_ejemplo_05_mutacion.png}{%
    En esta muestra, sigma = 1.0 alcanza la colina global 0 veces de 2000 y sigma = 3.0 la alcanza 110; cero eventos observados no prueba probabilidad cero}

% ================= primer_ejemplo_06_el_ciclo_completo =================
\section{El algoritmo genético completo}

\begin{frame}{El algoritmo genético completo}
inicializar -> [ SELECCIONAR -> CRUZAR -> MUTAR -> reemplazar ] x N
\end{frame}

\begin{frame}[fragile]{(1) el ciclo generacional}
\lstinputlisting[basicstyle=\ttfamily\fontsize{6.5pt}{7.8pt}\selectfont,firstline=138,lastline=162]{../src/primer_ejemplo_06_el_ciclo_completo.py}
\end{frame}

\begin{frame}[fragile]{(2) la gráfica de convergencia}
\lstinputlisting[basicstyle=\ttfamily\fontsize{7pt}{8.5pt}\selectfont,firstline=170,lastline=180]{../src/primer_ejemplo_06_el_ciclo_completo.py}
\end{frame}

\marcoresultado{El algoritmo genético completo}{../figuras/primer_ejemplo_06_el_ciclo_completo.png}{%
    Diez generaciones encuentran \texttt{x = +1.372, f = +0.706}, cerca de la cima más alta (referencia de 400 puntos: \texttt{x = +1.378, f = +0.706})}

% ================= primer_ejemplo_07_optimo_local =================
\section{Cuando falla}

\begin{frame}{Cuando falla}
SEMILLA = 52  ->  SEMILLA = 16
\end{frame}

\begin{frame}[fragile]{(1) SEMILLA}
\lstinputlisting[basicstyle=\ttfamily\fontsize{8pt}{9.6pt}\selectfont,firstline=32,lastline=32]{../src/primer_ejemplo_07_optimo_local.py}
\end{frame}

\marcoresultado{Cuando falla}{../figuras/primer_ejemplo_07_optimo_local.png}{%
    El mismo algoritmo se asienta en \texttt{x = -4.417, f = +0.073} y nunca sale}

\section{Siguiente}

\begin{frame}{Siguiente}
La semilla 52 encuentra la cima más alta observada; la 16 se queda en \(x = -4.417\), \(f = +0.073\) y no sale. Ninguna corrida prueba optimalidad global.

\vspace{0.8em}
\textbf{Lección 02 --- El flujo del algoritmo genético}

El ciclo que acabamos de escribir sigue siendo una ristra de comentarios sin nombre. La 02 nombra las cinco fases. Misma semilla, mismas cifras que el paso 6.
\end{frame}
\end{document}
